%%%%%%%%%%  scan idx 441  |  folio 428  %%%%%%%%%%
ment. On peut alors construire de proche en proche une résolution partielle graduée de \(\bigoplus_{n\ge0}J_0^n\), par des \(\cS(E_0)\)-modules gradués libres de type fini, de longueur supérieure à la dimension cohomologique de \(A_0\) ; le dernier noyau est alors un \(\cS(E_0)\)-module gradué, projectif de type fini en tant que \(\cS(E_0)\)-module non gradué, donc localement libre de type fini sur \(\Spec(A_0)\) en tant que \(\cS(E_0)\)-module gradué d'après VI 1.5.1. On obtient alors la résolution finie annoncée en remplaçant le dernier module de la résolution partielle par le noyau correspondant.

Si on tensorise la résolution ainsi construite par \(A\), on obtient alors un complexe gradué de \(\cS(E)\)-modules gradués localement libres de type fini, et il s'agit de voir que c'est une résolution de \(\bigoplus_{n\ge0}J^n\). Or le \(i\)-ième Module d'homologie du complexe ainsi obtenu est un \(\cS(E)\)-Module gradué dont la \(n\)-ième composante est \(\Tor_i^{A_0}(A,J_0^n)\), qui est nul pour \(i\ge1\) d'après (1.2.4). Par suite, d'après (1.2.3), on a construit ainsi une résolution du type voulu.

\textbf{Proposition 1.9.} Soient \(i:Y\to X\) une immersion fermée régulière, \(X'\) le \(X\)-schéma obtenu en faisant éclater \(Y\), et \(f:X'\to X\) le morphisme structural. Alors \(i\) et \(f\) sont des morphismes parfaits. (III 4.1.).

La propriété étant locale sur \(X\), on peut supposer donnés un \(\cO_X\)-Module localement libre de type fini \(E\), et un homomorphisme régulier \(u:E\to\cO_X\), ayant pour image l'idéal \(J\) de \(Y\). Le complexe de Koszul \(K_\bullet(u)\) est alors une résolution finie de \(\cO_Y\) par des \(\cO_X\)-Modules localement libres de type fini ; par conséquent, \(\cO_Y\) est un \(\cO_X\)-Module parfait. Il en résulte que toute immersion régulière est un morphisme parfait. D'après 1.8, \(f\) peut s'écrire comme le composé d'une immersion régulière \(X'\to\bfP(E)\) et du morphisme structural \(\bfP(E)\to X\) qui est lisse, donc parfait ; c'est donc
